\Needspace{60mm}
\subsubsection{结果分析}\label{sec:q3-4}


\textbf{（1）实验设置与评价指标。}本地实验采用本地模拟器生成的400个场景，每场含10至16个干扰源，位置在半径1800 m的圆形场地内按面积均匀分布。

设第 $i$ 场共有 $N_i$ 个干扰源，完成任务的总耗时为 $T_i$。评价时先将每场耗时除以该场干扰源数，再将各场结果取算术平均，同时报告平均每局耗时：
\begin{equation}\label{eq:q3-17}
\overline q=\frac1n\sum_{i=1}^{n}\frac{T_i}{N_i},
\qquad
\overline T=\frac1n\sum_{i=1}^{n}T_i.
\end{equation}
其中，$\overline q$ 的单位为秒/源，$\overline T$ 的单位为秒/局。总耗时包括清除最后一个干扰源后，为确认没有遗漏而进行的检测与移动，也包括光学精确定位失败时的耗时；此外记录完成全部清除的局数和实际清除的干扰源数。$\overline q$ 先对每场单独计算，再取平均，与直接将全部时间除以全部干扰源数的结果不同。

